% Part: first-order-logic
% Chapter: introduction
% Section: sentences

\documentclass[../../../include/open-logic-section]{subfiles}

\begin{document}

\olfileid{fol}{int}{snt}

% SEGMENT OLP-0144-S01
\olsection{\usetoken{P}{sentence}}

இப்போது !!{structure}s மற்றும் !!{formula}s க்கான நிறைவுறுத்தலின் (``இதில்
மெய்'') ஒரு (வரைபடமான) வரையறை நமக்குக் கிடைத்துள்ளது. ஆனால் அதற்கு இந்தக்
கூடுதல் கூறு---!!a{variable} ஒதுக்கீடு---தேவைப்படுகிறது; நமக்கு வேண்டியது
!!{sentence}s இன் வரையறை. ஒதுக்கீடுகளை எவ்வாறு நீக்குவது, !!{sentence}s
எவை?

முறையிய தருக்கத்திற்கான உங்கள் முதல் அறிமுகத்தில் !!{variable}s மற்றும்
அளவையடைகளுக்கு இடையிலான உறவைப் பற்றிய விவாதம் உங்களுக்கு நினைவிருக்கலாம்.
ஒரு அளவையடையைத் தொடர்ந்து எப்போதும் !!a{variable} வரும்; அந்த அளவையடை
பொருந்தும் !!{sentence} இன் பகுதியில் (அதன் ``வீச்சில்''), அந்த !!{variable}
அளவையடையால் ``கட்டுண்டது'' என்று புரிந்துகொள்கிறோம். !!{formula}s இல்
ஒவ்வொரு !!{variable} க்கும் இணையான அளவையடை இருக்க வேண்டும் என்று
கட்டாயமில்லை; இணையான அளவையடை இல்லாத !!{variable}s ``கட்டற்ற'' அல்லது
``கட்டுண்டிராத'' மாறிகள். கட்டற்ற !!{variable}s இல்லாத எல்லா !!{formula}s ஐயும்
!!{sentence}s எனக் கொள்வோம்.

மீண்டும், !!a{formula}~$!A$ இல் !!a{variable} ஒன்றின் நிகழ்வு எப்போது
கட்டுண்டது, எந்த அளவையடை அதை கட்டுகிறது, எப்போது கட்டற்றது என்பதற்கான
உள்ளுணர்வுக் கருத்து கடினமல்ல. இதைச் சோதிக்க நீங்கள் ஒரு முறையைக் கற்றிருக்கலாம்;
அதில் அடைப்புக்குறிகளை எண்ணுவதும் இருக்கலாம். இருந்தாலும் துல்லியமான
வரையறையை வலியுறுத்த வேண்டும். மேலும் !!{formula}s ஐத் தொகுத்தறிதலால்
வரையறுத்திருப்பதால், !!a{formula}~$!A$ இல் !!a{variable}~$x$ இன் கட்டற்ற மற்றும்
கட்டுண்ட நிகழ்வுகளையும் தொகுத்தறிதலால் வரையறுக்கலாம். நமது எளிய மொழிக்கு இது
பின்வருமாறு அமையலாம்:
\begin{enumerate}
  \item $!A$ அணுவாய்பாடாக இருந்தால், அதிலுள்ள~$x$ இன் எல்லா நிகழ்வுகளும்
  கட்டற்றவை (அதாவது,~$\Atom{\Obj P}{x}$ இல்~$x$ இன் நிகழ்வு கட்டற்றது).
  \item $!A$ ஆனது~$\lnot !B$ வடிவில் இருந்தால், $\lnot !B$ இல்~$x$ இன்
  நிகழ்வு கட்டற்றது அப்போதும் அப்போது மட்டுமே, அதற்குரிய~$x$ இன் நிகழ்வு
  $!B$ இல் கட்டற்றதாகும் (அதாவது,~$!B$ இல் உள்ள மாறிகளின் கட்டற்ற நிகழ்வுகள்
  எவையோ அவற்றுக்குரிய நிகழ்வுகளே~$\lnot !B$ இல் உள்ளன).
  \item $!A$ ஆனது $(!B \land !C)$ வடிவில் இருந்தால்,
  $(!B \land !C)$ இல்~$x$ இன் நிகழ்வு கட்டற்றது அப்போதும் அப்போது மட்டுமே,
  அதற்குரிய~$x$ இன் நிகழ்வு~$!B$ இல் அல்லது~$!C$ இல் கட்டற்றதாகும்.
  \item $!A$ ஆனது~$\lexists[x][!B]$ வடிவில் இருந்தால்,~$!A$ இல்~$x$ இன்
  எந்த நிகழ்வும் கட்டற்றதல்ல. அது~$\lexists[y][!B]$ வடிவில் இருந்து, $y$
  என்பது~$x$ இலிருந்து வேறுபட்ட !!{variable} எனில், $\lexists[y][!B]$ இல்
  $x$ இன் நிகழ்வு கட்டற்றது அப்போதும் அப்போது மட்டுமே, அதற்குரிய~$x$ இன்
  நிகழ்வு~$!B$ இல் கட்டற்றதாகும்.
\end{enumerate}

கட்டற்ற மற்றும் கட்டுண்ட நிகழ்வுகளுக்கான துல்லியமான வரையறை கிடைத்ததும்,
கட்டற்ற !!{variable}s இன் நிகழ்வுகள் இல்லாத !!{formula} எதுவும் !!a{sentence}
என்று எளிதாகக் கூறலாம்.

\end{document}
